% !TEX root = ~/Code/TN_LOWE_RECONSTRUCTION/main.tex
\section{Low-Energy Event Reconstruction\label{sec:RECONSTRUCTION}}
% Delete the text and write your Method(s) here:
%------------------
As a next step, after the generation, propagation and detector simulation of the \gls{mc} workflow, the reconstruction (\emph{reco}) step for these studies includes basic hit-finding algorithm and clustering schemes that applies both to the \gls{tpc} and \gls{pds}. This \emph{reco} stage typically is subdivided into two:

\begin{itemize}
    \item \emph{Reco1}: Conceptually, this part of the process is designed to provide the common building blocks for subsequent reconstruction stages: the \texttt{hits}. These are localized charge or photon depositions in the \gls{tpc} or \gls{pds}, respectively. As quanta of information, the algorithms that produce them are optimized for efficiency and accuracy by a full team of \gls{larsoft} developers. The \texttt{hits} used in this work are the standard \gls{larsoft} objects produced from the main reconstruction workflow and common to all \gls{dune} studies. Nevertheless, higher-level clustering algorithms have been specifically developed to group these \texttt{hits} into more complex objects that can be used for low-energy event reconstruction.
    \begin{itemize}
        
        \item \textbf{Common Preprocessing}: Both \gls{tpc} and \gls{pds} waveforms undergo a deconvolution step to remove the electronics response and noise spectrum. This process enhances the signal quality and prepares the waveforms for subsequent hit finding.
        
        \begin{itemize}
            \item \textbf{TPC Hit Finding}: 
            The \gls{roi}-trimmed waverforms are fitted to Gaussian Bells and integrated over time windows to form localized \texttt{hits}. These which are stored with time, position, and charge-profile information and form the fundamental building blocks for subsequent reconstruction stages.
            \item \textbf{TPC Disambiguation}: This step is specifically required for \gls{hd}, where induction wires turn around the central collection plane. This design allows for the reduction of the number of readout channels and simplifies, from a technical point of view, the extraction of the signals. However, it introduces an ambiguity in the hit position reconstruction, as a single readout channel can correspond to multiple physical wire positions (left of right of the collection plane). The disambiguation algorithm uses timing and charge information from all three planes to resolve these ambiguities and assign the correct wire section to each \texttt{hit}.
            \item \textbf{PDS Deconvolution}: The \gls{pds} waveforms are deconvolved to remove the electronics response and noise spectrum, similar to the \gls{tpc} deconvolution step. This process enhances the signal quality and prepares the waveforms for subsequent hit finding.
            \item \textbf{PDS Hit Finding}: The \gls{pds} deconvolved waveforms are processed to extract \emph{optical hits} or \texttt{ophits} (integrated photon pulses) with associated time, channel and \gls{pe} information. Instead of a simple Gaussian-fitting, a sliding-window algorithm is used to identify local maxima in the waveform and collect as much of the correlated scintillation light as possible into single \texttt{ophits}. This approach improves the timing and charge resolution of the detected scintillation signals.
        \end{itemize}
        
        \item \textbf{Specific Low-Level Clustering}: After the \texttt{hit} finding step, specific clustering algorithms have been developed to group these low-level objects into more complex structures that can be used for low-energy event reconstruction.
        \begin{itemize}
            \item \textbf{TPC Clustering}: The \gls{tpc} \texttt{hits} are grouped ``per-plane'' according to tight space and time constrains to form \texttt{clusters} (see \cref{fig:TPC_CLUSTERING}), which aim to combine all the charge from localized ionisation signals of single particles. This operation is performed for each of the \gls{tpc} planes. The clustering parameters found to maximize completeness while ensuring purity in the context of background are \num{3} adjacent channels and \SI{12.5}{\micro\second} time window (1~tick $=$ \SI{12.5}{\micro\second}). The resulting \texttt{per-plane clusters} contain the weighted average time and wire-position as well as the total charge.
            \item \textbf{TPC Matching}: The \texttt{per-plane clusters} from the three \gls{tpc} planes are matched together starting from the collection plane and associating the corresponding induction plane clusters based on the shortest time (TPC-plane clocks are synchronized so that the drift time between planes is accounted for). The resulting \texttt{matched clusters} contain information from all three planes and most importantly, a reconstructed 2.5D position $<t,y,z>$.
        \end{itemize}
    \end{itemize}

    \item \emph{Reco2}: This second part of the reconstruction process focuses on higher-level object formation and the association between \gls{tpc} and \gls{pds} systems. The following operations are performed:
    \begin{itemize}
        \item \textbf{PDS Flash Finding}: The \texttt{ophits} are grouped following a simpler clustering logic with a much broader space range to account for the larger spread of scintillation light. Contrary to the \gls{tpc} clustering scheme, the algorithm implemented for the \gls{pds} relies on fixed time and distance parameters that collect \texttt{adjacent ophit} surrounding a \texttt{trigger ophit} with the highest \gls{pe} count. This algorithm contemplates the possibility of collecting into more than one flash the same adjacent hits. This feature can be turned on or off depending on the analysis needs. The resulting \texttt{flashes} contain time, position and total \gls{pe} information for subsequent matching with \gls{tpc} objects (see \cref{fig:FLASH_FINDING}).
        
        \item \textbf{Low-Energy Event Finding}: The \emph{Reco1} \gls{tpc} \texttt{cluster} are further grouped into \texttt{lowe events} containing one \texttt{primary} and any number of found \texttt{adjacent clusters} as well as a candidate list of matching \gls{pds} \texttt{flashes}:
        \begin{itemize}
            \item \textbf{TPC Adjacent Cluster Finding}: Each \texttt{cluster} is scanned for smaller \texttt{adjacent clusters} within a radial distance of \SI{1}{\meter} (see \cref{fig:ADJ_CLUSTER_FINDING}). The resulting \texttt{primary cluster} (the highest charge cluster within \SI{1}{\meter}) is the fundamental component to form the \texttt{lowe event}.
            \item \textbf{TPC-PDS Flash Matching}: Each \texttt{primary cluster} is assigned to a list of flashes via flash-matching algorithm which takes into account the reconstructed vertices of the clusters and the time and position of the flashes (find more details to the matching algorithm in~\cref{sec:FLASH_MATCHING}).
        \end{itemize}
    \end{itemize}
\end{itemize}

The clustering algorithms developed for this purpose are detailed in~\cref{sec:CLUSTERING_TPC,sec:CLUSTERING_PDS} and can be found in the following repositories: \href{https://github.com/DUNE/duneopdet/tree/develop/duneopdet/SolarNuUtils}{duneopdet} \& \href{https://github.com/DUNE/dunereco/tree/develop/dunereco/LowEUtils}{dunereco}; as well as the associated producer modules: \href{https://github.com/DUNE/duneopdet/blob/develop/duneopdet/OpticalDetector/SolarOpFlash_module.cc}{SolarOpFlash\_module.cc} \& \href{https://github.com/DUNE/dunereco/blob/develop/dunereco/LowEUtils/SolarCluster_module.cc}{SolarCluster\_module.cc}. A custom solar workflow leading to the production of these reconstructed objects has been included in the standard production workflow files:

\begin{itemize}
    \item \href{https://github.com/DUNE/dunesw/blob/develop/fcl/dunefd/reco/solar_reco_dune10kt_1x2x6.fcl}{solar\_reco\_dune10kt\_1x2x6.fcl}
    \item \href{https://github.com/DUNE/dunesw/blob/develop/fcl/dunefdvd/reco/solar_reco_dunevd10kt_1x8x14_3view_30deg.fcl}{solar\_reco\_dunevd10kt\_1x8x14\_3view\_30deg.fcl}.
\end{itemize}

The calibration and characterization of the \gls{cc} neutrino signal are based on the true information of the main electron of the interaction. If not otherwise mentioned, the plots in this section are produced from a MARLEY generated \gls{cc} neutrino sample with central \gls{apa} background. The first step of the reconstruction process has already happened when clustering the hits of the \gls{tpc} and \gls{pds} detector subsystems in the simulation's \emph{reco1} stage (see \cref{sec:EVENT_DEFINITION}). Next, we evaluate the detection efficiency and completeness of the collection- induction-matched \gls{tpc} clusters (see \cref{sec:DETECTION_EFFICIENCY}). Subsequently, the event electron is corrected in charge for the drift attenuation before computing a hit-dependent calibration factor that recover the particle's original kinetic energy (see \cref{sec:ELECTRON_CALIBRATION}). Finally, the average calibration factor is applied to all adjacent clusters of the event candidate and the total neutrino energy is computed from four different algorithms (see \cref{sec:NEUTRINO_CALIBRATION}). During the \gls{tpc} calibration process, the drift coordinate is assumed from truth. A discussion on how the intrinsic background contamination of the detector affects the signal reconstruction at these low energies and the effects on drift coordinate resolution (as a result of \gls{tpc}-\gls{pds} match) are presented in~\cref{sec:RESULTS}.

\subsection{Flash Matching\label{sec:FLASH_MATCHING}}
The first step in the signal reconstruction process is the attempt to match the reconstructed primary \ac{TPC} cluster with its corresponding \ac{PDS} flash. Note that this step is not strictly required for the energy reconstruction, as the drift attenuation correction can be applied statistically. However, the matching is essential for the correct reconstruction of the neutrino position and, most importantly, background fiducialization. This is particularly important for low--energy solar neutrino analyses, where the signal rate is low and the high energy tails of the background spectra can leak into the signal region. The matching is also important for supernova neutrino analyses, where the time structure of the neutrino burst provides valuable information about the supernova dynamics.

The matching is performed by evaluating the distance between the \ac{TPC} vertex and the PDS flash position, as well as the time difference between the \ac{TPC} cluster and the PDS flash. The matching efficiency is defined as the fraction of \ac{TPC} clusters that have a corresponding pure PDS flash. \cref{fig:FLASH_MATCHING} shows the flash--matching efficiency as a function of the distance to the PDS plane and the neutrino energy. The matching efficiency is high for distances up to \SI{1}{\meter} and decreases for larger distances. The matching efficiency is also higher for higher neutrino energies. This is expected, as the PDS flashes are more likely to be detected for higher energy depositions.

\begin{figure}[ht]
    \centering
    \includegraphics[width=\linewidth]{Images/hd_1x2x6_lateralAPA_marley_official_TPC_Cluster_Efficiency_Scan.png}
    \caption[width=1\linewidth]{Flash--matching efficiency. (left) As a function of distance to the PDS plane. (right) As a function of neutrino energy.\label{fig:FLASH_MATCHING}}
\end{figure}

For any given \ac{TPC} cluster, the matching strategy follows these steps:

\begin{enumerate}
    \item\textbf{Distance Calculation}: Compute the distance between the \ac{TPC} cluster vertex and all PDS flash positions.
    \item\textbf{Time Difference Calculation}: Calculate the time difference between the \ac{TPC} cluster and the PDS flash. This must be smaller than the drift time of the electrons in the LAr.
    \item\textbf{Candidate Selection Criteria}: A PDS flash is considered adjacent to a \ac{TPC} flash if the distance $(\text{TPC}_\text{y} - \text{PDS}_\text{y})^2 / \text{AdjOpFlashY}^2 + (\text{TPC}_\text{Z} - \text{PDS}_\text{Z})^2 / \text{AdjOpFlashZ}$ is smaller than 1, where \( \text{AdjOpFlashY} \) and \( \text{AdjOpFlashZ} \) are the adjusted matching distances in the 2 projected coordinates, respectively.
    \item\textbf{Matching Criteria}: The algorithm consider 3 additional cuts to prevent the random matching of background cluster:
    \begin{itemize}
        \item The minimum number of hits in the PDS flash: AdjOpFlashMinNHitCut (by default 3).
        \item The minimum number of PEs in the PDS flash: AdjOpFlashMinPECut (by default 20).
        \item The maximum ratio of AdjOpFlashMaxPE / AdjOpFlashPE: AdjOpFlashMaxPERatioCut (by default 1). This is intended to prevent the matching of background PDS flashes, wich are less homogeneous in the number of PEs compared to the signal PDS flashes.
    \end{itemize}
\end{enumerate}

Select the PDS flash with the highest number of PE as the matching candidate for the \ac{TPC} cluster. If no PDS flash is found within the matching criteria, the \ac{TPC} cluster is considered unmatched.

From this matching one can reconstruct the drift coordinate of the \ac{TPC} cluster, which is essential for the correct reconstruction of the neutrino energy. The drift coordinate is computed from the time between the \ac{TPC} vertex and the PDS flash by taking into account the drift velocity of the electrons in the \ac{LAr}.
\subsection{\acrshort{pds} Clustering}\label{sec:CLUSTERING_PDS}
Similarly to the \gls{tpc} clustering, the \gls{pds} signals are also clustered in order to reconstruct the flashes. The flash finding scheme is illustrated in~\cref{fig:FLASH_FINDING}. The algorithm takes into account the time and channel proximity of the \texttt{ophits} to cluster them into flashes. Each flash is then assigned a vertex and a time, which are used for the flash-matching with the \gls{tpc} clusters. The clustering is performed independently for each \gls{pds} plane in the \gls{vd} geometry, while in the \gls{hd} geometry all \texttt{ophits} are clustered together as there is only one \gls{pds} plane. The resulting flash contains information about the total number of \gls{pe}s, the weighted average time of the flash, and the flash vertex (computed as the weighted average position of all \gls{xa} channels). The choice of the time and channel proximity thresholds is crucial, as it determines the size and purity of the flashes. A smaller threshold will result in more, smaller flashes, while a larger threshold will produce fewer, larger flashes. The optimal thresholds are determined based on the expected signal characteristics and the background conditions of the detector. In this work, we found that a time window of \SI{100}{\micro\second} and a radial distance of \SI{1}{\meter} provides a good balance between flash completeness and purity in the context of background. The resulting flashes are then used for the flash-matching with the \gls{tpc} clusters, which is described in the next section.

\begin{figure}[ht!]
    \centering
    \includegraphics[width=0.8\linewidth]{Images/Schematics/FLASH_FINDING_SCHEME.png}
    \caption{Flash finding scheme. Optical channels (\gls{xa}s) are represented as black squares. The darkness of their filling represents the number of \gls{pe}s detected per \texttt{hit}. The color represents the \texttt{ophit}'s precedence: signal (red) and background (blue). The dashed circle shows the \texttt{ophits} grouped into a single flash. Notice how the algorithm allows for the same \texttt{ophit} to be part of more than one flash.\label{fig:FLASH_FINDING}}
\end{figure}

The flashes are produced from \texttt{ophits} by clustering these \gls{pds} signals in time and channel proximity. All \texttt{ophits} within a given time window (\texttt{OpFlashAlgoMinTime} - \texttt{OpFlashAlgoMaxTime}) and distance (\texttt{OpFlashAlgoRad}) are grouped together around a single trigger \texttt{ophit}, which is defined to be the one with the highest number of \gls{pe}s. The clustering is performed independently for each \gls{pds} plane in the \gls{vd} detector. The resulting flash contains information about the total number of \gls{pe}s, the weighted average time of the flash, and the flash vertex(computed as the weighted average position of all \gls{xa} channels).

Once the set of flashes is formed, the matching is performed by evaluating the distance between the \gls{tpc} vertex and the \gls{pds} flash, as well as their time difference. Flashes can be formed from signal or background activity. For any given flash, there is a non--negligible chance clustering pileup (where \texttt{ophits} formed from background activity get clustered into flashes where the trigger \texttt{ophit} originates from signal) or channel pileup (where \texttt{ophits} themselves are formed from a mixture of simultaneously arriving signal and background photons) occur. Therefore, we distinguish between:

\begin{itemize}
    \item \textbf{Signal flashes}: those that contain at least some \texttt{ophit} originating from the signal activity. In any realistic scenario, these flashes will be contaminated with at least some amount of photons originating from background activity. Signal flashes can be subdivided into:
    \begin{itemize}
        \item \textbf{High Purity}: those that contain mostly \texttt{ophits} originating from the signal activity. Their reconstruction is likely to provide an accurate estimate of the true neutrino interaction time and good vertex reconstruction.
        \item \textbf{Low Purity}: those that contain mostly \texttt{ophits} originating from background activity. Their reconstruction is likely to provide a smeared estimate of the true neutrino interaction time and poor vertex reconstruction.     
    \end{itemize}
    \item \textbf{Background flashes}: those that contain only \texttt{ophits} originating from background activity. These flashes can be mistakenly matched to a \gls{tpc} cluster, leading to a wrong drift coordinate reconstruction. As these are not correlated in time with the neutrino interaction, they will lead to a random and flat drift coordinate assignment within the detector active volume.
\end{itemize}

This distinction will become important as we will see in~\cref{sec:VERTEX_RESOLUTION}. Before attempting any full reconstruction based on these matching, it is important to understand the reach of the matching algorithm.

\begin{figure}[ht]
    \centering
    \includegraphics[width=.49\linewidth]{Images/vd_1x8x14_3view_30deg_marley_official_x_scan.png}
    \includegraphics[width=.49\linewidth]{Images/vd_1x8x14_3view_30deg_marley_official_plane_scan.png}
    \caption[width=1\linewidth]{Flash--matching efficiency for the \gls{vd} geometry. Results are shown as a function of distance to the \gls{pds} plane (left) and \gls{pds}--plane match (right).\label{fig:FLASH_MATCHING_VD}}
\end{figure}
\subsection{Electron--Lifetime Correction}\label{sec:ELECTRON_LIFETIME}
\ac{TPC} hits are first grouped into spatial clusters using a dedicated clustering algorithm. Each event contains a \textbf{main cluster}, corresponding to the primary outgoing electron, and zero or more \textbf{adjacent clusters}, interpreted as gamma-induced ``blips''. The first step in calibration involves correcting the total ionization charge of each cluster for the electron--lifetime attenuation in the \ac{LAr} (see~\cref{eq:ELECTRON_LIFETIME}).

\begin{figure}[ht!]
    \centering
    \includegraphics[width=0.49\linewidth]{Images/hd_1x2x6_centralAPA_marley_SignalParticleKineticEnergy_vs_NHits.png}
    \includegraphics[width=0.49\linewidth]{Images/hd_1x2x6_centralAPA_marley_Charge_Correction_Hist2D.png}
    \caption[width=1\linewidth]{(left) Neutrino energy correlation with number of hits in the primary cluster. (right) Drift attenuation for the main electron cluster.\label{fig:DRIFT_ATTENUATION}}
\end{figure}

\begin{equation}
    \dfrac{Q_\text{raw}(n_\text{hit})}{E^\text{truth}_e} = F_0(n_\text{hit}) \cdot \exp\left(\frac{t^\text{drift}}{\tau_e}\right),
\end{equation}\label{eq:ELECTRON_LIFETIME}

where \( Q_{\text{raw}} \) is the observed (uncorrected) cluster charge, \( E^\text{truth}_e \) is the true main electron energy of the interaction, \( F_0 \) is the charge--to--energy correction factor, \( t^\text{drift} \) is the true drift time of the electron, and \( \tau_e \) is the electron lifetime measured here independently and cross-checked against the simulation input.

\begin{figure}[ht!]
    \centering
    \includegraphics[width=\linewidth]{Images/hd_1x2x6_centralAPA_marley_Charge_Corrected_Charge.png}
    \caption[width=1\linewidth]{Drift attenuation correction.\label{fig:CHARGE_CORRECTION}}
\end{figure}

Assuming the main cluster corresponds to a fully-contained electron, we take its known true energy (from simulation) as a reference and compute the charge--to--energy conversion factor for each $(n_\text{hit})$ threshold. Additionally, to assess the impact of this assumption, with the goal of characterizing signal clusters in the context of background which limits the clustering parameters, one can also investigate a ``cheated charge''. This is computed for all clusters, adding the surrounding charge that has been backtracked to the main electron of the interaction. A calibration based on this cheated clustering gives us an understanding of the losses that arise from a limited clustering ability in the context of background.

\begin{equation}
    E_{\text{cluster}} = \alpha(n_{\text{hit}}, E_{\text{true}}) \cdot Q_{\text{corr}}.
\end{equation}

After the truth drift correction has been applied to the data, the charge correction factor is immediately obtained.

\begin{figure}[ht!]
    \centering
    \includegraphics[width=.49\linewidth]{Images/hd_1x2x6_centralAPA_marley_official_Charge_Correction_Factor.png}
    \includegraphics[width=.49\linewidth]{Images/hd_1x2x6_centralAPA_marley_official_ElectronCharge_Correction_Factor.png}
    \caption{Per hit correction factor distributions. (left) Primary cluster charge. (right) Cheated electron cluster charge.\label{fig:CHARGE_CORRECTION_FACTOR}}
\end{figure}

\subsection{Electron Energy Calibration}\label{sec:ELECTRON_CALIBRATION}
Even after lifetime and charge calibration, small systematic biases may remain in the reconstructed electron energy due to recombination modelling, imperfect response, or clustering inefficiencies. We apply a linear correction to account for this residual offset:

\begin{equation}
    E_{e}^{\text{final}} = \left(E_{\text{cluster}} - n\right) / m,
\end{equation}

where \( m \) and \( n \) are derived from the linear fits to the simulated \( E_{\text{true}} \) versus \( E_{\text{reco}} \) distributions as shown in \cref{fig:CHARGE_CALIBRATION}. The calibration is performed separately for the main electron cluster and the cheated electron cluster. The calibration factor is computed as a function of the number of hits in the cluster (see~\cref{fig:ENERGY_CALIBRATION}), which serves as a proxy for the electron energy.

\begin{figure}[ht!]
    \centering
    \includegraphics[width=\linewidth]{Images/hd_1x2x6_centralAPA_marley_Charge_Calibration_2D_NHits1.png}
    \caption[width=1\linewidth]{Charge calibration for the main electron cluster. (left) Corrected charge over true electron energy. (center) Linear fit to reco energy over true electron energy. (rigth) Comparison between (true - reco) for corrected (before linear fit calibration) and calibrated (after linear fit calibration) reco energy distributions.\label{fig:CHARGE_CALIBRATION}}
\end{figure}

\begin{figure}[ht!]
    \centering
    \includegraphics[width=.49\linewidth]{Images/hd_1x2x6_centralAPA_marley_Charge_Calibration.png}
    \includegraphics[width=.49\linewidth]{Images/hd_1x2x6_centralAPA_marley_ElectronCharge_Calibration.png}
    \caption{Per hit calibration results. (left) Primary cluster charge. (right) Cheated electron cluster charge. Blue shows the slope and red the intercept values of the linear fits.\label{fig:ENERGY_CALIBRATION}}
\end{figure}